\section{Introduction}

This document describes a \textbf{Counter} atomic model in both UA-DEVS (Uncertainty-Aware DEVS) and IA-DEVS (Interval-Approximated DEVS) formalisms. The Counter model accumulates incoming events and outputs the accumulated count when reset.

\section{Informal Description}

The Counter model receives two types of events:
\begin{itemize}
    \item \textbf{Add}: Increments an internal counter by 1.
    \item \textbf{Reset}: Outputs the current accumulated count and resets the counter to 0.
\end{itemize}

After processing an \textit{Add} event, the model becomes passive (waits indefinitely for the next event). After receiving a \textit{Reset} event, the model outputs the accumulated count with time advance of 0 (immediate output) and then resets its internal counter to 0.

\section{UA-DEVS Counter Specification}

Following the UA-DEVS formalism from ``Uncertainty Aware Discrete-Event System Simulation''~\cite{vicino2021uncertainty}, we define the Counter model as:

\subsection{Model Definition}

The Counter atomic model is defined as the tuple:
\[
Counter_{UA} = \langle X, Y, S, \delta_{int}, \delta_{ext}, \lambda, ta \rangle
\]
where:

\begin{itemize}
    \item $X = \mathcal{P}(\{\text{Add}, \text{Reset}\}) \setminus \emptyset$ is the set of input values (non-empty subsets of event types). This allows representing uncertainty about which event(s) occurred, including:
    \begin{itemize}
        \item $\{\text{Add}\}$: definitely an Add event
        \item $\{\text{Reset}\}$: definitely a Reset event
        \item $\{\text{Add}, \text{Reset}\}$: uncertain which event occurred (could be either)
    \end{itemize}
    \item $Y = \mathcal{P}(\mathbb{N}) \setminus \emptyset$ is the set of output values (non-empty subsets of non-negative integers). This allows outputting sets of possible count values to represent uncertainty.
    \item $S = \mathcal{P}(\{\text{passive}, \text{output}\} \times \mathbb{N}) \setminus \emptyset$ is the state space (non-empty subsets of phase-count pairs). This captures uncertainty about the current state. Any valid state in $S$ can serve as an initial state when instantiated by the simulator.
\end{itemize}

\subsection{Transition Functions}

	extbf{Internal Transition ($\delta_{int}$):}
	extbf{Internal Transition ($\delta_{int}$):}
\[
\delta_{int}(S') = \{(\text{passive}, 0) \mid (\text{output}, c) \in S'\}
\]
    extbf{External Transition ($\delta_{ext}$):}

Given a state set $S' \subseteq S$, elapsed time set $E \subseteq \mathbb{R}^+_0$, and input set $X' \subseteq X$:
\[
\delta_{ext}(S', E, X') = \bigcup_{s \in S', e \in E, x \in X'} \delta_{ext}^{single}(s, e, x)
\]
where $\delta_{ext}^{single}$ operates on individual elements:
\[
\delta_{ext}^{single}((\phi, c), e, x) = 
\begin{cases}
    \{(\text{passive}, c + 1)\} & \text{if } x = \text{Add} \\
    \{(\text{output}, c)\} & \text{if } x = \text{Reset}
\end{cases}
\]

The elapsed time $e$ does not affect the state for this model. When $X' = \{\text{Add}, \text{Reset}\}$, the result includes both the incremented state and the output state, capturing the uncertainty.

\bigskip

	extbf{Output Function ($\lambda$):}
\[
\lambda(S') = \{c \mid (\text{output}, c) \in S'\}
\]
Output the set of all possible count values from states in the output phase.

\bigskip

\textbf{Time Advance ($ta$):}
\[
ta(S') = \begin{cases}
    \{\infty\} & \text{if } \forall (\phi, c) \in S', \phi = \text{passive} \\
    \{0\} & \text{if } \forall (\phi, c) \in S', \phi = \text{output} \\
    \{0, \infty\} & \text{if } \exists (\text{passive}, c_1), (\text{output}, c_2) \in S'
\end{cases}
\]
The time advance returns a set representing all possible next event times. When states are mixed between passive and output phases, both immediate (0) and infinite time advances are possible.

\subsection{Behavior Description}

The Counter model can be initialized by the simulator with any valid state from $S$. For typical usage, it might start at $\{(\text{passive}, 0)\}$ and wait indefinitely. Upon receiving input $\{\text{Add}\}$, it transitions to $\{(\text{passive}, 1)\}$ and remains passive. Upon receiving input $\{\text{Reset}\}$, it transitions to $\{(\text{output}, c)\}$ with $ta = \{0\}$. Immediately, the output function $\lambda$ emits the set $\{c\}$, and then the internal transition $\delta_{int}$ resets the state to $\{(\text{passive}, 0)\}$.

When there is uncertainty in the input (e.g., $\{\text{Add}, \text{Reset}\}$), the external transition produces multiple possible states. For example, from state $\{(\text{passive}, 3)\}$ with input $\{\text{Add}, \text{Reset}\}$, the result is $\{(\text{passive}, 4), (\text{output}, 3)\}$, capturing both possibilities.

\section{IA-DEVS Counter Specification}

In IA-DEVS, we work with intervals to approximate sets of values to work with uncertainty in computable context. All components (states, times, inputs, outputs) are represented as intervals $[l, u]$ where $l$ and $u$ are the lower and upper bounds.

\subsection{Order Relations}

For intervals to be well-defined in IA-DEVS, we must establish total order relations on the base types:

\begin{itemize}
    \item \textbf{Input event order}: $\text{Add} < \text{Reset}$. This ordering allows representing input uncertainty as intervals. For example, $[\text{Add}, \text{Reset}]$ represents uncertainty about which event occurred.
    
    \item \textbf{Phase order}: $\text{passive} < \text{output}$. The phases are ordered with passive coming first.
    
    \item \textbf{State order}: States $(\phi, c) \in \{\text{passive}, \text{output}\} \times \mathbb{N}$ are ordered lexicographically, first by phase, then by count value:
    \[
    (\phi_1, c_1) < (\phi_2, c_2) \iff \phi_1 < \phi_2 \text{ or } (\phi_1 = \phi_2 \text{ and } c_1 < c_2)
    \]
    This means all passive states come before all output states, and within each phase, states are ordered by their count value.
    
    \item \textbf{Output order}: Outputs are natural numbers with the standard ordering $\mathbb{N}$.
\end{itemize}

These orderings enable interval arithmetic and ensure that interval operations (union, intersection, etc.) produce well-defined results.

\subsection{Model Definition}

The IA-DEVS Counter model is:
\[
Counter_{IA} = \langle X_I, Y_I, S_I, \Delta_{int}, \Delta_{ext}, \Lambda, TA \rangle
\]

where:

\begin{itemize}
    \item $X_I = \mathcal{I}(\{\text{Add}, \text{Reset}\})$ is the set of input intervals over the ordered set of event types. Typical intervals include $[\text{Add}, \text{Add}]$, $[\text{Reset}, \text{Reset}]$, or $[\text{Add}, \text{Reset}]$ (representing uncertainty).
    
    \item $Y_I = \mathcal{I}(\mathbb{Z}_{[0, \text{MAX\_INT}]})$ is the set of output intervals, where $\mathbb{Z}_{[0, \text{MAX\_INT}]}$ represents bounded integers from 0 to MAX\_INT as implemented by the computational platform. Intervals may be right-open, e.g., $[c, +\infty)$, when the upper bound reaches or exceeds MAX\_INT.
    
    \item $S_I = \mathcal{I}(\{\text{passive}, \text{output}\} \times \mathbb{Z}_{[0, \text{MAX\_INT}]})$ is the state interval space using bounded integers. State intervals can be:
    \begin{itemize}
    \item \textbf{Closed}: $[(\phi_l, c_l), (\phi_u, c_u)]$ when $c_u < \text{MAX\_INT}$ or $\phi_u = \text{output}$ prevents further growth
    \item \textbf{Right-open}: $[(\phi_l, c_l), +\infty)$ when count growth is unbounded or reaches MAX\_INT
    \end{itemize}
    
    Examples:
    \begin{itemize}
    \item $[(\text{passive}, 0), (\text{passive}, 5)]$: count in [0, 5], definitely passive (closed)
    \item $[(\text{output}, 3), (\text{output}, 7)]$: count in [3, 7], definitely output (closed)
    \item $[(\text{passive}, 2), (\text{output}, 3)]$: mixed phase and count uncertainty (closed)
    \item $[(\text{passive}, 100), +\infty)$: count $\geq 100$, passive phase (right-open, unbounded growth)
    \item $[(\text{passive}, \text{MAX\_INT}), +\infty)$: count at or beyond MAX\_INT (right-open, saturation)
    \end{itemize}
\end{itemize}

Note: The initial state is determined by the simulator when instantiating the model, not by the model definition itself.

\subsection{Computational Constraints}

In the IA-DEVS implementation, the count values are represented using fixed-size integer types (e.g., \texttt{int32\_t}, \texttt{int64\_t}). This introduces the following considerations:

\begin{itemize}
    \item When $c_u = \text{MAX\_INT}$ and an Add event arrives, the resulting state interval must be right-open: $[(\phi, c_l + 1), +\infty)$ to represent the possibility of overflow or saturation.
    
    \item The implementation may choose to either:
    \begin{enumerate}
        \item \textbf{Saturate}: clamp values at MAX\_INT (conservative approximation)
        \item \textbf{Signal overflow}: use right-open intervals to indicate unbounded growth
    \end{enumerate}
    
    \item Right-open intervals preserve the soundness of the approximation: they over-approximate the set of reachable states, ensuring no valid trajectory is excluded.
\end{itemize}

\subsection{Approximated Transition Functions}

\textbf{Internal Transition ($\Delta_{int}$):}

The internal transition handles states in the output phase, resetting them to passive:

\[
\Delta_{int}([(\phi_l, c_l), (\phi_u, c_u)]) = 
\begin{cases}
    [(\text{passive}, 0), (\text{passive}, 0)] & \text{if } \phi_l = \phi_u = \text{output} \\
    [(\text{passive}, 0), (\text{passive}, c_u)] & \text{if } \phi_l = \text{passive}, \phi_u = \text{output}
\end{cases}
\]

\textbf{Explanation:}
\begin{itemize}
    \item When all states are in output phase ($[\text{output}, \text{output}]$), all are reset to $(\text{passive}, 0)$.
    \item When phases are mixed ($[\text{passive}, \text{output}]$), the bound in passive state was never in output phase and does not transition; it is now the value of the upper bound, while the lower bound resets its count to 0.
\end{itemize}

\bigskip

\textbf{External Transition ($\Delta_{ext}$):}

Given state interval $[(\phi_l, c_l), (\phi_u, c_u)]$, elapsed time interval $[e_l, e_u]$, and input interval $[x_l, x_u]$:

\begin{itemize}
    \item If $[x_l, x_u] = [\text{Add}, \text{Add}]$:
    \[
    \Delta_{ext}([(\phi_l, c_l), (\phi_u, c_u)], [e_l, e_u], [\text{Add}, \text{Add}]) = [(\text{passive}, c_l + 1), (\text{passive}, c_u + 1)]
    \]
    If the upper bound is right-open (unbounded), the result remains right-open: $[(\text{passive}, c_l + 1), +\infty)$.
    
    \item If $[x_l, x_u] = [\text{Reset}, \text{Reset}]$:
    \[
    \Delta_{ext}([(\phi_l, c_l), (\phi_u, c_u)], [e_l, e_u], [\text{Reset}, \text{Reset}]) = [(\text{output}, c_l), (\text{output}, c_u)]
    \]
    The count values are preserved; only the phase changes to output. If the interval was right-open, it remains right-open.
    
    \item If $[x_l, x_u] = [\text{Add}, \text{Reset}]$ (uncertain input):
    \[
    \Delta_{ext}([(\phi_l, c_l), (\phi_u, c_u)], [e_l, e_u], [\text{Add}, \text{Reset}]) = [(\text{passive}, c_l), (\text{output}, c_u)]
    \]
    This captures uncertainty: lower bound corresponds to Add (increment to passive), upper bound to Reset (change to output).
\end{itemize}

The elapsed time $e$ does not affect the state in this model, as the counter only responds to discrete events.
Note that technically, the input interval could also be open, but all open intervals in a 2-element set can be transformed into a closed one.

\bigskip

\textbf{Output Function ($\Lambda$):}

The output function extracts count values from states in the output phase, preserving interval closure properties:

For a state interval $I_s = [ (\phi_l, c_l), (\phi_u, c_u) ]$ with lower closure flag $lc$ and upper closure flag $uc$:

\[
\Lambda(I_s) = I_c
\]

where $I_c$ is the count interval with:
\begin{itemize}
    \item Lower bound: $c_l$ if $\phi_l = \text{output}$, otherwise the smallest count in output phase
    \item Upper bound: $c_u$ if $\phi_u = \text{output}$, or $\infty+$ if right unbounded
    \item Lower closure: matches $lc$ 
    \item Upper closure: matches $uc$
\end{itemize}

\textbf{Examples:}
\begin{itemize}
    \item $[(\text{output}, 3), (\text{output}, 7)]$ with both closed $\Rightarrow$ output $[3, 7]$ closed
    \item $(\text{output}, 3), (\text{output}, 7)]$ with left open, right closed $\Rightarrow$ output $(3, 7]$
    \item $[(\text{output}, 3), (\text{output}, 7))$ with left closed, right open $\Rightarrow$ output $[3, 7)$
    \item $[(\text{output}, 3), +\infty)$ right-open $\Rightarrow$ output $[3, +\infty)$
    \item $[(\text{passive}, 3), (\text{output}, 7)]$ mixed phases $\Rightarrow$ output includes counts from output states [3, 7] ( $(\text{output}, 3)$ is greater than $(\text{passive}, 3)$ )
\end{itemize}

The key principle: the output interval's closure properties (open/closed at lower/upper bounds) directly mirror the closure properties of the input state interval, ensuring that boundary inclusion/exclusion semantics are preserved.

\bigskip

\textbf{Time Advance ($TA$):}
\[
TA([(\phi_l, c_l), (\phi_u, c_u)]) = 
\begin{cases}
    \emptyset & \text{if } \phi_l = \phi_u = \text{passive} \\
    [0, 0] & \text{if } \phi_l = \phi_u = \text{output} \\
    [0, +\infty) & \text{if } \phi_l = \text{passive}, \phi_u = \text{output}
\end{cases}
\]

In IA-DEVS, an empty time advance interval $\emptyset$ represents passivity (infinite time advance). A $[0, 0]$ interval means immediate scheduling. When phases are mixed, the result is $[0, +\infty)$ (right-open) representing both immediate (for output states) and infinite (for passive states) time advances.

\subsection{Implementation Notes}

For the C++ implementation in the Cadmia library:

\begin{itemize}
    \item Use \texttt{cadmia::modeling::decimal<3>} for time representation (millisecond precision).
    \item Represent states as intervals: \texttt{cadmia::modeling::interval<state\_t>} where \texttt{state\_t} contains phase enum and integer count (phase, count).
    \item Input types: use an enumeration or integer encoding for \textit{Add} and \textit{Reset}.
    \item Follow the existing pattern from \texttt{processor.hpp} and \texttt{generator.hpp}: define \texttt{state\_t}, \texttt{time\_t}, \texttt{input\_t}, \texttt{output\_t}, and their interval counterparts.
    \item Implement static functions: \texttt{internal\_transition}, \texttt{external\_transition}, \texttt{output}, and \texttt{time\_advance}.
    \item Use \texttt{interval::empty\_interval()} for passive (infinite) time advance.
    \item Validate inputs and throw \texttt{std::invalid\_argument} for invalid states or transitions.
\end{itemize}

\subsection{Example Scenario}

	extbf{Initial State:} $[(\text{passive}, 0), (\text{passive}, 0)]$, $TA = \emptyset$ (passive).

\textbf{Step 1:} External input $[\text{Add}, \text{Add}]$ arrives.
\[
	ext{New state: } [(\text{passive}, 1), (\text{passive}, 1)], \quad TA = \emptyset
\]

\textbf{Step 2:} External input $[\text{Add}, \text{Add}]$ arrives.
\[
	ext{New state: } [(\text{passive}, 2), (\text{passive}, 2)], \quad TA = \emptyset
\]

\textbf{Step 3:} External input $[\text{Reset}, \text{Reset}]$ arrives.
\[
	ext{New state: } [(\text{output}, 2), (\text{output}, 2)], \quad TA = [0, 0]
\]

\textbf{Step 4:} Internal transition (output $[2, 2]$ and reset).
\[
\Lambda = [2, 2], \quad \text{New state: } [(\text{passive}, 0), (\text{passive}, 0)], \quad TA = \emptyset
\]

The cycle repeats for further events.

\section{Conclusion}

The Counter model illustrates a simple yet complete IA-DEVS atomic model with two input event types and interval-based state representation. This specification provides the formal foundation for implementing the model in the Cadmia C++ library and writing comprehensive unit tests.

\begin{thebibliography}{1}
\bibitem{vicino2021uncertainty}
Damián Vicino, Gabriel Wainer, and Olivier Dalle.
\textit{Uncertainty on Discrete-Event System Simulation}.
ACM Transactions on Modeling and Computer Simulation (TOMACS), 2021.
\end{thebibliography}
